\section*{Chapitre \ref{ch:ml:large-language-models-in-finance} --- Grands modèles de langage en finance}

\begin{solution}{exo:ml:large-language-models-in-finance:1}
Seuls les mois de juillet à décembre de la deuxième année, après la date de coupure, permettent de mesurer le jugement, et seulement si la date de coupure du fournisseur est
exacte et que le modèle n'a pas été mis à jour depuis. Les dix-huit mois qui la précèdent mesurent un mélange de jugement et de
mémoire de ce qui a suivi chaque titre.
\end{solution}

\begin{solution}{exo:ml:large-language-models-in-finance:2}
$\text{idf} = \log(1 + (540 - 45 + 0.5)/(45 + 0.5)) = \log 11.89 = 2.48$. À longueur moyenne, le normalisateur vaut $k_1 = 1.5$, donc
le facteur de terme vaut $1\times2.5/(1 + 1.5) = 1$ et la contribution vaut 2.48.
\end{solution}

\begin{solution}{exo:ml:large-language-models-in-finance:3}
Prime (entre avant et après)~: modèle contaminé $64.8 - 59.2 = 5.6$ points, modèle propre $56.2 - 56.3 = -0.1$. Perte par anonymisation
entre les dates de coupure~: modèle contaminé $64.8 - 57.0 = 7.8$, modèle propre $56.2 - 56.0 = 0.2$.
\end{solution}

\begin{solution}{exo:ml:large-language-models-in-finance:4}
Son score ne dépend que des mots pour lesquels il dispose de plongements, des synonymes de la métrique et des mots de contexte~; les noms d'entreprises et les années
ne figurent pas dans son vocabulaire, si bien que les 135 passages portant sur la bonne métrique sont à égalité et que celui en tête concerne la bonne entreprise par
hasard (1.1~\%). Ajouté à \omterm{def:ml:large-language-models-in-finance:rag}{BM25}, il remonte de la même façon chaque passage sur la bonne métrique, y compris le passage de l'autre
entreprise qui partage le synonyme de la question et que \omterm{def:ml:large-language-models-in-finance:rag}{BM25} classe déjà en tête~: les deux scores ne divergent que là où la recherche dense est
aveugle.
\end{solution}

\begin{solution}{exo:ml:large-language-models-in-finance:5}
Il a été entraîné sur un tiers de titres en plus (7\,988 jusqu'au jour 1\,210 au lieu de 5\,972 jusqu'au jour 907), donc il estime mieux les effets des événements, et plus
récemment. Après les deux dates de coupure, il n'a vu aucune des réponses~: le gain est donc de la généralisation, au sens du
chapitre, un jugement acquis grâce aux données. L'intérêt du protocole est que cette comparaison est équitable et que celle entre les dates de coupure ne l'est pas.
\end{solution}

\begin{solution}{exo:ml:large-language-models-in-finance:6}
Chaque titre précède la date de coupure du modèle, si bien que les 70~\% mêlent mémoire et jugement~; la comparaison avec un modèle tf-idf
entraîné et testé hors échantillon est inéquitable. Retirer les noms des entreprises ne constitue pas une anonymisation complète (dates, produits, personnes
et des formulations uniques permettent d'identifier l'événement), et les 65~\% après cette date montrent seulement qu'une part de la mémoire a subsisté ou que le texte est
informatif, sans dire lequel. Évaluer le score sur les titres postérieurs à la date de coupure, anonymiser les noms et les dates, et comparer avec les références sur les
mêmes titres.
\end{solution}

\begin{solution}{exo:ml:large-language-models-in-finance:7}
Modifier le dernier jeton laisse inchangés les logits à toutes les positions antérieures (le test d'acceptation le vérifie à $10^{-6}$).
Le masque d'attention est triangulaire supérieur, avec $-\infty$ au-dessus de la diagonale, si bien qu'après le softmax la position $t$ accorde un poids nul
aux positions situées après $t$~; les couches feed-forward agissent position par position, donc rien d'autre ne mélange les positions.
\end{solution}

\begin{solution}{exo:ml:large-language-models-in-finance:8}
Question~: les signaux dont le coefficient d'information a le plus baissé entre le trimestre précédent et le dernier, avec la
mesure de décroissance du livre~7 (chapitre~13). Outils~: le registre des signaux, le magasin de données point-in-time dont la date as-of est fixée à la
fin de trimestre, le calculateur d'IC, un outil de tracé~; aucun outil ne peut écrire. Contrôles~: le garde-fou as-of, le nombre minimal
d'observations par signal, un ajustement pour tests multiples lorsqu'on classe de nombreux signaux. Journal~: la tâche, la version du modèle, le \omterm{def:ml:large-language-models-in-finance:prompt}{prompt},
chaque \omterm{def:ml:large-language-models-in-finance:agent}{appel d'outil} avec ses arguments et les empreintes de ses résultats, les refus, la réponse finale et le code qui la recalcule sans
l'agent.
\end{solution}

\begin{solution}{pb:ml:large-language-models-in-finance:1}
\textbf{Partie I.}
\begin{enumerate}
\item Phrases formées d'un jeton de semaine, de l'entreprise, du titre et du résultat après une flèche, pour chaque titre antérieur à sa
date de coupure~; interrogé en soumettant un \omterm{def:ml:large-language-models-in-finance:prompt}{prompt} jusqu'à la flèche et en comparant les probabilités de «~hausse~» et de «~baisse~».
\item 106\,681 paramètres et 505 jetons. Le fait qu'un modèle mémorise dépend de la fréquence à laquelle il a vu un document et des clés qui
l'identifient, pas seulement de la taille~; le protocole (découpage par période, anonymisation) mesure la mémoire dans tout modèle.
\item 61.6~\%~: le signe de l'effet injecté est juste dans 66.0~\% des cas sur les 1\,462 titres comportant un effet, et les
autres 550 sont des tirages à pile ou face. Le bruit de la réaction (3~\%) est important par rapport aux effets.
\item Pour que les substituts soient des mots que le modèle a déjà vus~; sinon le \omterm{def:ml:large-language-models-in-finance:memo}{test d'anonymisation} mesurerait la réaction du modèle
à des jetons inconnus, et non la perte de sa mémoire.
\end{enumerate}
\textbf{Partie II.}
\begin{enumerate}[resume]
\item Modèle propre 69.7~\%, modèle contaminé 65.4~\%~; versions anonymisées 57.7~\% et 57.2~\%. Les deux ont mémorisé les résultats de leurs titres d'entraînement
titres, indexés par semaine et par entreprise~; sans ces clés, ils retombent sur leur jugement.
\item Modèle contaminé 64.8~\%, modèle propre 56.2~\%~: le modèle contaminé a vu ces titres et leurs résultats.
\item Il ramène le modèle contaminé à 57.0~\% (une perte de 7.8 points) et le modèle propre à 56.0~\% (0.2).
\item Lopez-Lira et Tang ont évalué sur des titres postérieurs à la date de coupure de GPT-4~; Glasserman et Lin ont constaté que les titres anonymisés
obtenaient de meilleurs résultats dans l'échantillon, car la connaissance générale de l'entreprise distrayait le modèle davantage que l'anticipation ne l'aidait.
\end{enumerate}
\textbf{Partie III.}
\begin{enumerate}[resume]
\item Même mot~: \omterm{def:ml:large-language-models-in-finance:rag}{BM25} 100~\%, dense 0~\%, hybride 100~\%, \omterm{def:ml:large-language-models-in-finance:rag}{BM25} enrichi 86.3~\%. Autre mot~: 5.0~\%, 1.1~\%, 18.2~\%, 89.1~\%.
\item Il refuse 11.9~\% des questions auxquelles on peut répondre, précisément celles dont le passage en tête désigne une autre entreprise ou une autre année et
auraient reçu une mauvaise réponse~; il s'abstient sur les 300 questions sans réponse, qui recevraient sinon toutes un nombre.
\item Il répond à 80.8~\% des questions as-of à partir d'un dépôt publié plus tard, et a raison dans 19.2~\% des cas~; avec le
garde-fou, 100~\%.
\item Chaque \omterm{def:ml:large-language-models-in-finance:agent}{appel d'outil} avec ses arguments, la date as-of, le résultat ou le refus, ainsi que la version du modèle et le \omterm{def:ml:large-language-models-in-finance:prompt}{prompt}.
\end{enumerate}
\textbf{Partie IV.}
\begin{enumerate}[resume]
\item \emph{Le modèle qui savait.} Prime de \omterm{def:ml:large-language-models-in-finance:memo}{mémorisation} de 5.6 points pour le modèle contaminé et $-0.1$ pour le modèle propre~;
perte par anonymisation de 7.8 points et 0.2.
\item Demander sa date de coupure et l'historique de ses versions, les résultats par période par rapport à la date de coupure, une exécution anonymisée, et un essai en direct ou
postérieur à la date de coupure contre une référence par dictionnaire.
\item Extraction, étiquetage et synthèse avec la source citée, et rédaction de code avec tests~: des sorties qu'une personne ou un test
peut vérifier.
\item Les informations privilégiées, les données clients et tout ce que la firme ne publierait pas.
\item Il lance plus d'expériences, plus vite, sans mémoire de ce qu'il a essayé, et exploite toutes les données renvoyées par les outils~; sans un
garde-fou et un journal, il surapprend et utilise de l'information future à grande échelle.
\item Figer la version du modèle, le \omterm{def:ml:large-language-models-in-finance:prompt}{prompt}, les réglages de décodage et l'index de recherche~; journaliser chaque appel~; conserver un chemin de code qui recalcule
le résultat sans le modèle.
\item Lorsque la tâche est stable, que des données étiquetées existent et que le \omterm{def:ml:large-language-models-in-finance:prompt}{prompt} laisse de la précision ou des coûts inexploités~; et lorsque le modèle
peut être hébergé là où les données peuvent être transmises.
\item Avant sa date de coupure, mémoire et jugement se mêlent~; après elle, seul le jugement.
\end{enumerate}
\end{solution}

\begin{solution}{iq:ml:large-language-models-in-finance:1}
La date du dernier texte des données d'entraînement. Avant cette date, le modèle a pu lire l'événement et ses suites, et un
backtest sur cette période comporte un biais d'anticipation qu'aucune hygiène des données côté société ne peut supprimer.
\iqlookfor{biais d'anticipation dans les données d'entraînement du modèle, et la parade (évaluation postérieure à la date de coupure).}
\end{solution}

\begin{solution}{iq:ml:large-language-models-in-finance:2}
Rechercher des passages, les placer dans le \omterm{def:ml:large-language-models-in-finance:prompt}{prompt}, puis générer une réponse qui s'appuie sur eux. Évaluer la recherche sur des
paires question--passage étiquetées (rappel en $k$), et la génération à partir du bon passage (fidélité, correspondance exacte), afin que
les échecs soient imputés à la bonne étape.
\iqlookfor{les deux étapes et une métrique distincte pour chacune.}
\end{solution}

\begin{solution}{iq:ml:large-language-models-in-finance:3}
Comparer l'exactitude avant et après la date de coupure annoncée~; anonymiser les noms et les dates et mesurer la perte~; tester sur des
événements que le modèle ne peut pas connaître (postérieurs à la coupure, ou synthétiques)~; vérifier s'il sait compléter mot pour mot le texte du titre.
\iqlookfor{le découpage par période et l'anonymisation, idéalement un test de complétion.}
\end{solution}

\begin{solution}{iq:ml:large-language-models-in-finance:4}
\omterm{def:ml:large-language-models-in-finance:rag}{BM25} pour les noms, les tickers, les nombres et les termes exacts~; des plongements pour les paraphrases~; en pratique les deux, le plongement servant à
enrichir la requête ou reclasser les candidats de \omterm{def:ml:large-language-models-in-finance:rag}{BM25}, mesuré sur des questions étiquetées tirées des documents propres au desk.
\iqlookfor{ce que chacun manque, et une hybridation mesurée.}
\end{solution}

\begin{solution}{iq:ml:large-language-models-in-finance:5}
Placer la date de référence dans le harnais, et non dans le \omterm{def:ml:large-language-models-in-finance:prompt}{prompt}~: chaque outil de données filtre selon l'heure de publication et refuse les
résultats postérieurs, et chaque appel est journalisé. Le prouver par le journal et par un test qui injecte des documents futurs et vérifie qu'ils sont refusés.
\iqlookfor{l'application dans les outils, un journal, et un test avec données futures insérées.}
\end{solution}

\begin{solution}{iq:ml:large-language-models-in-finance:6}
Le modèle et la version utilisés, ainsi que leurs dates de coupure~; les résultats découpés avant et après la coupure~; un backtest anonymisé~; le
suivi en direct depuis le lancement~; les horodatages des scores par rapport aux actualités~; les coûts et la capacité~; et un essai gratuit
sur des données postérieures à la date de coupure.
\iqlookfor{la découpe selon la coupure, l'anonymisation, l'historique en direct et les horodatages.}
\end{solution}
